%========================================
% MatinBook — Stage 03: Layout Test
% Version: v1.1
% Purpose:
%   Validate the layout and navigation systems of MatinBook v1.1:
%     1. Page geometry (frontmatter/mainmatter split)
%     2. Header/footer (twoside)
%     3. Heading styles (Boyer/Stewart scale)
%     4. Table of contents
%     5. Footnotes
%     6. Line spacing
%     7. Margin notes (RTL)
%     8. Page styles (plain, fancy)
%
% Compilation:
%   xelatex -shell-escape stage03-layout.tex
%   xelatex -shell-escape stage03-layout.tex
%   (2 passes required for cover + cross-references)
%
% Expected outcome:
%   A professionally typeset PDF demonstrating that every
%   layout system of MatinBook works correctly, including the
%   wide margin for RTL margin notes.
%========================================

%------------------------------------------------------------------------
% Search Path (for tests directory)
%------------------------------------------------------------------------
\makeatletter
\def\input@path{{../}}
\makeatother

%------------------------------------------------------------------------
% Document Class
%------------------------------------------------------------------------
\documentclass{matinbook}

%------------------------------------------------------------------------
% Book Metadata
%------------------------------------------------------------------------
\title{آزمون صفحه‌آرایی و ساختار}
\author{تیم توسعه MatinBook}
\date{مهر ۱۴۰۵}

\booktitle{آزمون صفحه‌آرایی و ساختار}

\renewcommand{\repository}{github.com/matinmahmoudi-cs/matinbook}

%========================================================================
% DOCUMENT
%========================================================================
\begin{document}

%------------------------------------------------------------------------
% FRONT COVER
%------------------------------------------------------------------------
\makecover
    {آزمون صفحه‌آرایی و ساختار}
    {ارزیابی هندسه، سربرگ و عناوین در MatinBook}
    {تیم توسعه MatinBook}
    {مهر ۱۴۰۵}

%------------------------------------------------------------------------
% FRONT MATTER (symmetric margins)
%------------------------------------------------------------------------
\frontmatter
\frontmattergeometry

% Title page
\maketitle

% Copyright / Colophon
\thispagestyle{empty}
\vfill
\begin{center}
    {\large\textbf{شناسنامه آزمون}}

    \vspace{0.8cm}

    \begin{tabular}{rl}
        عنوان: & آزمون صفحه‌آرایی و ساختار \\
        نسخه: & \matinbookversion \\
        تاریخ: & \matinbookdate \\
        موتور: & \lr{XeLaTeX} \\
        مجوز: & \lr{MIT} \\
    \end{tabular}

    \vspace{0.8cm}

    {\small این سند بخشی از مجموعه آزمون‌های یکپارچگی \lr{MatinBook} است.}
\end{center}
\clearpage

% Preface
\chapter*{پیشگفتار}
\addcontentsline{toc}{chapter}{پیشگفتار}

این سند، سومین آزمون از مجموعه‌ی آزمون‌های یکپارچگی \lr{MatinBook}
نسخه‌ی \lr{\matinbookversion} است. هدف آن، اعتبارسنجی کامل سیستم
صفحه‌آرایی و ساختار سند در هشت حوزه‌ی متمایز است: هندسه‌ی صفحه،
سربرگ و پابرگ، استایل عناوین، فهرست مطالب، پانویس‌ها، فاصله‌ی
خطوط، یادداشت‌های حاشیه‌ای و استایل‌های صفحه.

در یک کتاب فنی فارسی، صفحه‌آرایی حرفه‌ای تعیین‌کننده‌ی کیفیت
نهایی اثر است. حاشیه‌های دقیق، سربرگ‌های خوانا، عناوین
سلسله‌مراتبی و فهرست مطالب کارآمد، همگی به خواننده کمک می‌کنند
تا محتوا را سریع‌تر درک کند و مسیر خود را در کتاب گم نکند.

هر بخش از این آزمون، با مثال‌های عملی و ارجاعات متقابل همراه
است تا خواننده بتواند درستی رفتار هر زیرسیستم را در بافت
واقعی یک کتاب ارزیابی کند.

\clearpage

% Table of Contents
\tableofcontents
\clearpage

% List of Figures
\listoffigures
\clearpage

% List of Tables
\listoftables
\clearpage

%------------------------------------------------------------------------
% MAIN MATTER (wide margin for notes)
%------------------------------------------------------------------------
\mainmatter
\mainmattergeometry

%========================================================================
% CHAPTER 1 — PAGE GEOMETRY
%========================================================================
\chapter{هندسه‌ی صفحه}
\label{chap:geometry}

\section{فلسفه‌ی هندسه در MatinBook}
\label{sec:geometry-philosophy}

\lr{MatinBook} از یک هندسه‌ی صفحه‌ی \textbf{دوطرفه} با حاشیه‌ی
بیرونی پهن استفاده می‌کند. این طراحی، از کتاب‌های درسی
دانشگاهی مانند \lr{Boyer} و \lr{Stewart} الهام گرفته شده
است که در آن‌ها حاشیه‌ی بیرونی، فضای کافی برای یادداشت‌های
توضیحی، شکل‌های کوچک و فرمول‌های مکمل فراهم می‌کند.

صفحه در \lr{MatinBook} با ابعاد \pnum{۱۷.۸} × \pnum{۲۵.۴}
سانتی‌متر (نسبت طلایی نزدیک به \lr{7 × 10} اینچ) تنظیم شده
است. این ابعاد، استاندارد نشر آکادمیک بین‌المللی است و
خوانایی بالایی در متن‌های فنی دارد.

\mnote{ابعاد \pnum{۱۷.۸} × \pnum{۲۵.۴} سانتی‌متر، نسبت طلایی نزدیک به \lr{7 × 10} اینچ است.}

\section{حاشیه‌های متن اصلی}
\label{sec:geometry-mainmatter}

در متن اصلی (\lr{mainmatter})، هندسه‌ی صفحه به این صورت است:

\begin{itemize}
    \item \textbf{حاشیه‌ی داخلی:} \pnum{۱.۵} سانتی‌متر (نزدیک عطف)
    \item \textbf{حاشیه‌ی بیرونی:} \pnum{۵.۹} سانتی‌متر (پهن، برای یادداشت)
    \item \textbf{عرض ناحیه‌ی یادداشت:} \pnum{۳.۵} سانتی‌متر
    \item \textbf{فاصله‌ی یادداشت از متن:} \pnum{۰.۶} سانتی‌متر
    \item \textbf{حاشیه‌ی بالا و پایین:} \pnum{۱.۷} سانتی‌متر
    \item \textbf{فضای صحافی:} \pnum{۰.۴} سانتی‌متر
\end{itemize}

عرض متن اصلی در این حالت، حدود \pnum{۱۰.۳} سانتی‌متر است.
این عرض، برای متن فارسی با فونت \lr{XB Niloofar} در اندازه‌ی
\pnum{۱۰pt} بهینه است و طول خطوط را در محدوده‌ی خوانای
\pnum{۶۰}--\pnum۷۵ نویسه نگه می‌دارد.

\section{حاشیه‌های پیش‌متن و پس‌متن}
\label{sec:geometry-frontmatter}

در پیش‌متن (\lr{frontmatter}) و پس‌متن (\lr{backmatter})،
هندسه‌ی صفحه به حالت متقارن تغییر می‌کند. در این حالت:

\begin{itemize}
    \item \textbf{حاشیه‌ی داخلی:} \pnum{۱.۵} سانتی‌متر
    \item \textbf{حاشیه‌ی بیرونی:} \pnum{۱.۵} سانتی‌متر (متقارن)
    \item \textbf{ناحیه‌ی یادداشت:} حذف می‌شود
    \item \textbf{عرض متن:} حدود \pnum{۱۴.۳} سانتی‌متر
\end{itemize}

این تغییر هندسه، با دو دستور \lr{\textbackslash frontmattergeometry}
و \lr{\textbackslash mainmattergeometry} انجام می‌شود. دلیل
این تفکیک، آن است که پیش‌متن (حق نشر، پیشگفتار، فهرست مطالب)
و پس‌متن (منابع، نمایه) به یادداشت حاشیه‌ای نیاز ندارند و
می‌توانند از عرض کامل متن استفاده کنند.

\section{تست متن طولانی}
\label{sec:geometry-longtext}

این یک پاراگراف طولانی است که برای ارزیابی نحوه‌ی شکسته
شدن خطوط و پر شدن صفحه نوشته شده است. هدف ما این است که
ببینیم آیا متن به‌درستی در صفحه جای می‌گیرد و حاشیه‌ها
رعایت می‌شوند یا خیر. \lr{MatinBook} از تنظیمات پیشرفته‌ی
\lr{microtype} برای بهبود ترازبندی و جلوگیری از سرریز
خطوط استفاده می‌کند.

فضای اضافی برای صحافی (\lr{bindingoffset}) در نظر گرفته
شده است تا متن در قسمت داخلی صفحه (نزدیک عطف کتاب) گم
نشود و خوانایی کتاب در حالت باز حفظ گردد. این تنظیم،
در چاپ حرفه‌ای کتاب، امری ضروری است.

\section{تست فاصله‌ی پاراگراف‌ها}
\label{sec:geometry-paragraphs}

این پاراگراف اول است. فاصله بین پاراگراف‌ها باید مناسب باشد
--- نه آنقدر زیاد که متن از هم گسیخته به نظر برسد، و نه
آنقدر کم که پاراگراف‌ها در هم ادغام شوند.

این پاراگراف دوم است. تنظیمات پیش‌فرض \lr{MatinBook} شامل
\lr{\textbackslash parindent=1em} برای تورفتگی ابتدای
پاراگراف و \lr{\textbackslash parskip=0pt} برای فاصله‌ی
بین پاراگراف‌هاست. این مقادیر، استاندارد نشر فارسی است.

این پاراگراف سوم است. توجه کنید که تورفتگی فقط برای
پاراگراف‌های دوم به بعد اعمال می‌شود و پاراگراف اول هر
بخش تورفتگی ندارد. این رفتار، به‌طور خودکار توسط
\lr{MatinBook} مدیریت می‌شود.

\section{پانویس‌ها}
\label{sec:geometry-footnotes}

این متن دارای یک پانویس است\footnote{این یک پانویس آزمایشی
است که باید در پایین صفحه نمایش داده شود. فاصله پانویس‌ها
از متن و از یکدیگر باید مناسب باشد. خط جداکننده پانویس،
۳۰٪ عرض متن است.}.

و این هم یک پانویس دیگر\footnote{پانویس دوم برای تست فاصله
بین پانویس‌ها. \lr{\textbackslash footnotesep=8pt} تنظیم
شده است.} برای اطمینان از صحت تنظیمات.

%========================================================================
% CHAPTER 2 — HEADERS AND FOOTERS
%========================================================================
\chapter{سربرگ و پابرگ}
\label{chap:headers}

\section{فلسفه‌ی طراحی سربرگ}
\label{sec:headers-philosophy}

\lr{MatinBook} از بسته‌ی \lr{fancyhdr} برای مدیریت سربرگ‌ها
و پابرگ‌ها استفاده می‌کند. طراحی سربرگ‌ها به‌صورت \textbf{دوطرفه}
است تا در چاپ کتاب، اطلاعات در محل درست خود قرار گیرند.

در صفحات فرد (راست‌چین)، سربرگ در سمت راست قرار می‌گیرد و
شماره صفحه در چپ. در صفحات زوج (چپ‌چین)، ترتیب برعکس است.
این طراحی، استاندارد نشر حرفه‌ای است.

\section{محتوای سربرگ}
\label{sec:headers-content}

در \lr{MatinBook}، سربرگ صفحات فرد و زوج شامل نام فصل جاری
است. نام فصل، از دستور \lr{\textbackslash chaptermark} گرفته
می‌شود که به‌طور خودکار توسط کلاس به‌روزرسانی می‌شود.

\begin{itemize}
    \item \textbf{صفحات فرد:} نام فصل در بالا راست
    \item \textbf{صفحات زوج:} نام فصل در بالا چپ
    \item \textbf{صفحات شروع فصل:} فقط پابرگ (بدون سربرگ)
\end{itemize}

\mnote{صفحات شروع فصل، استایل \lr{plain} دارند و فقط شماره صفحه را در پابرگ نشان می‌دهند.}

\section{محتوای پابرگ}
\label{sec:headers-footer}

پابرگ در \lr{MatinBook} شامل شماره صفحه است. خط جداکننده‌ی
سربرگ با ضخامت \pnum{۰.۴pt} طراحی شده است، در حالی که
پابرگ بدون خط است. این طراحی مینیمال، به خواننده کمک می‌کند
بدون حواس‌پرتی روی محتوا تمرکز کند.

\section{تست چندین صفحه}
\label{sec:headers-multipage}

برای ارزیابی کامل سربرگ و پابرگ در چندین صفحه، این پاراگراف‌های
طولانی برای پر کردن صفحات نوشته شده‌اند. هدف آن است که ببینیم
آیا نام فصل به‌درستی در سربرگ همه‌ی صفحات به‌روزرسانی می‌شود
و شماره صفحه در محل درست قرار می‌گیرد.

\subsection{زیربخش اول}
\label{sec:headers-sub1}

محتوای علمی می‌تواند شامل توضیحات مفصل درباره یک موضوع باشد.
برای مثال، در این بخش می‌توانیم درباره اهمیت صفحه‌آرایی
در کتاب‌های فارسی صحبت کنیم. صفحه‌آرایی خوب باعث می‌شود
خواننده تجربه‌ی بهتری از مطالعه داشته باشد و مطالب را
با سرعت و دقت بیشتری درک کند.

\subsection{زیربخش دوم}
\label{sec:headers-sub2}

یکی از چالش‌های صفحه‌آرایی فارسی، ترکیب متن فارسی با
عناصر لاتین مانند کدهای برنامه‌نویسی، فرمول‌های ریاضی
و جداول است. \lr{MatinBook} با استفاده از \lr{xepersian}
و تنظیمات دقیق، این مشکل را به‌خوبی حل کرده است.

\subsection{زیربخش سوم}
\label{sec:headers-sub3}

فاصله بین بخش‌ها و زیربخش‌ها باید به گونه‌ای باشد که
سلسله‌مراتب مطالب به وضوح مشخص شود. عناوین اصلی باید
برجسته‌تر از عناوین فرعی باشند و خواننده بتواند
به‌راحتی ساختار مطلب را دنبال کند.

%========================================================================
% CHAPTER 3 — HEADING STYLES
%========================================================================
\chapter{استایل عناوین}
\label{chap:headings}

\section{فلسفه‌ی مقیاس عناوین}
\label{sec:headings-philosophy}

\lr{MatinBook} از مقیاس عناوین \textbf{Boyer/Stewart} استفاده
می‌کند که در کتاب‌های درسی دانشگاهی معتبر است. این مقیاس،
سلسله‌مراتب زیر را تعریف می‌کند:

\begin{itemize}
    \item \textbf{شماره‌ی فصل:} \pnum{۲۴pt}
    \item \textbf{عنوان فصل:} \pnum{۲۰pt}
    \item \textbf{بخش:} \pnum{۱۵pt}
    \item \textbf{زیربخش:} \pnum{۱۳pt}
    \item \textbf{زیرزیربخش:} \pnum{۱۱pt}
\end{itemize}

\section{عنوان فصل}
\label{sec:headings-chapter}

عنوان فصل با استایل خاص \lr{MatinBook} شامل یک شماره‌ی
بزرگ و کم‌رنگ در پس‌زمینه، عنوان با فونت درشت، و یک خط
رنگی در پایین است. این طراحی از کتاب‌های فنی مدرن الهام
گرفته شده است.

\section{عنوان بخش}
\label{sec:headings-section}

این یک \lr{\textbackslash section} است. عنوان بخش با
جعبه‌ی رنگی شماره و فونت \pnum{۱۵pt} نمایش داده می‌شود.
شماره‌ی بخش در یک \lr{colorbox} با پس‌زمینه‌ی آبی روشن
قرار می‌گیرد.

\subsection{عنوان زیربخش}
\label{sec:headings-subsection}

این یک \lr{\textbackslash subsection} با فونت \pnum{۱۳pt}
و شماره‌ی آبی کم‌رنگ است. زیربخش‌ها، سطح دوم سلسله‌مراتب
را تشکیل می‌دهند.

\subsubsection{عنوان زیرزیربخش}
\label{sec:headings-subsubsection}

این یک \lr{\textbackslash subsubsection} با فونت \pnum{۱۱pt}
و شماره‌ی آبی خیلی کم‌رنگ است. زیرزیربخش‌ها، عمیق‌ترین
سطح عنوان‌بندی در \lr{MatinBook} را تشکیل می‌دهند.

\section{تست سلسله‌مراتب کامل}
\label{sec:headings-hierarchy}

\subsection{سطح دوم: زیربخش}
\label{sec:headings-level2}

این متن زیر یک زیربخش قرار دارد. باید تورفتگی مناسبی
نسبت به عنوان بخش داشته باشد.

\subsubsection{سطح سوم: زیرزیربخش}
\label{sec:headings-level3}

این متن زیر یک زیرزیربخش قرار دارد. عمیق‌ترین سطح
عنوان‌بندی در \lr{MatinBook} است.

\subsection{زیربخش دیگر}
\label{sec:headings-other}

برای مقایسه‌ی فاصله‌ی بین دو زیربخش متوالی.

%========================================================================
% CHAPTER 4 — TABLE OF CONTENTS
%========================================================================
\chapter{فهرست مطالب}
\label{chap:toc}

\section{ساختار فهرست}
\label{sec:toc-structure}

فهرست مطالب \lr{MatinBook} با استفاده از بسته‌ی \lr{tocloft}
ساخته می‌شود و دارای ویژگی‌های زیر است:

\begin{itemize}
    \item \textbf{فونت فصل‌ها:} سیاه و پررنگ
    \item \textbf{شماره صفحه‌ی فصل‌ها:} آبی پررنگ
    \item \textbf{نقطه‌چین:} با رنگ آبی کم‌رنگ
    \item \textbf{تورفتگی:} برای سطوح مختلف
    \item \textbf{فاصله:} \pnum{۸pt} قبل از فصل، \pnum{۴pt} قبل از بخش
\end{itemize}

\section{فهرست‌های فرعی}
\label{sec:toc-sublists}

علاوه بر فهرست مطالب، \lr{MatinBook} از فهرست تصاویر
(\lr{\textbackslash listoffigures})، فهرست جداول
(\lr{\textbackslash listoftables}) و فهرست الگوریتم‌ها
(\lr{\textbackslash listofalgorithms}) نیز پشتیبانی می‌کند.

در این آزمون، فهرست تصاویر و فهرست جداول در ابتدای سند
نمایش داده شده‌اند. فهرست الگوریتم‌ها فقط در صورت استفاده
از محیط \lr{algorithm} فعال می‌شود.

\section{ارجاع متقابل}
\label{sec:toc-cref}

می‌توانیم به بخش‌های مختلف سند ارجاع دهیم. برای مثال،
در \cref{chap:summary} نتایج نهایی این آزمون را مشاهده
خواهید کرد. همچنین \cref{tab:layout-status} وضعیت
اجزای صفحه‌آرایی را نشان می‌دهد.

%========================================================================
% CHAPTER 5 — PAGE STYLES
%========================================================================
\chapter{استایل‌های صفحه}
\label{chap:pagestyles}

\section{صفحه‌ی شروع فصل}
\label{sec:pagestyles-plain}

صفحات شروع فصل، استایل \lr{plain} دارند و فاقد سربرگ
هستند. فقط پابرگ با شماره صفحه نمایش داده می‌شود. این
طراحی تمیز، توجه خواننده را به ابتدای فصل جلب می‌کند
و فضای بصری بیشتری برای عنوان فصل فراهم می‌آورد.

\section{صفحات معمولی}
\label{sec:pagestyles-fancy}

صفحات معمولی، استایل \lr{fancy} دارند و شامل سربرگ کامل
(نام فصل + شماره صفحه) و پابرگ هستند. خط جداکننده‌ی
سربرگ با ضخامت \pnum{۰.۴pt} طراحی شده است.

\section{فاصله‌ی خطوط}
\label{sec:pagestyles-spacing}

فاصله‌ی خطوط \lr{MatinBook} روی \pnum{۱.۱۵} تنظیم شده است
(\lr{\textbackslash setstretch\{1.15\}}). این مقدار، برای
متن فارسی با فونت \lr{XB Niloofar} بهینه است و خوانایی
را افزایش می‌دهد بدون اینکه فضای زیادی اشغال کند.

\mnote{فاصله‌ی خطوط \pnum{۱.۱۵} از سبک کتاب‌های درسی \lr{Boyer/Stewart} الهام گرفته شده است.}

%========================================================================
% CHAPTER 6 — MARGIN NOTES
%========================================================================
\chapter{یادداشت‌های حاشیه‌ای}
\label{chap:marginnote}

\section{فلسفه‌ی یادداشت حاشیه‌ای}
\label{sec:marginnote-philosophy}

یکی از ویژگی‌های متمایزکننده‌ی \lr{MatinBook}، پشتیبانی
از یادداشت‌های حاشیه‌ای در حاشیه‌ی بیرونی پهن است. این
یادداشت‌ها برای توضیحات مکمل، ارجاعات سریع و نکات کلیدی
به کار می‌روند.

\section{انواع یادداشت}
\label{sec:marginnote-types}

دو نوع یادداشت حاشیه‌ای در \lr{MatinBook} وجود دارد:

\begin{itemize}
    \item \lr{\textbackslash mnote\{...\}} — یادداشت متنی ساده
    \item \lr{\textbackslash marginfig\{file\}\{caption\}} — شکل با کپشن
\end{itemize}

هر دو نوع، به‌طور خودکار در حاشیه‌ی بیرونی صفحه (راست
در صفحات فرد، چپ در صفحات زوج) قرار می‌گیرند و جهت
متن آن‌ها راست‌به‌چپ است.

\section{محدودیت‌ها}
\label{sec:marginnote-limits}

یادداشت‌های حاشیه‌ای محدودیت‌هایی نیز دارند:

\begin{itemize}
    \item \textbf{داخل \lr{tcolorbox}:} قابل استفاده نیست
    \item \textbf{داخل \lr{figure}:} قابل استفاده نیست
    \item \textbf{نیاز به هندسه‌ی \lr{mainmatter}:} در \lr{frontmatter} ناحیه‌ی یادداشت حذف می‌شود
    \item \textbf{کامپایل چندباره:} برای موقعیت‌یابی پایدار، حداقل ۲ بار کامپایل لازم است
\end{itemize}

%========================================================================
% CHAPTER 7 — SUMMARY
%========================================================================
\chapter{خلاصه‌ی نتایج}
\label{chap:summary}

\section{وضعیت اجزای صفحه‌آرایی}
\label{sec:summary-layout}

جدول \cref{tab:layout-status} وضعیت اجزای مختلف
صفحه‌آرایی \lr{MatinBook} را نشان می‌دهد.

\begin{matintable}{وضعیت اجزای صفحه‌آرایی}{tab:layout-status}
\begin{tblr}{
    width = \textwidth,
    colspec = {l X[c] X[c]},
    hlines, vlines,
    colsep = 8pt,
    rowsep = 4pt,
    row{1} = {font=\bfseries},
}
بخش & وضعیت & توضیحات \\
هندسه‌ی صفحه & \textcolor{matingreen}{فعال} & \pnum{۱۷.۸} × \pnum{۲۵.۴}cm، حاشیه‌ی بیرونی \pnum{۵.۹}cm \\
سربرگ & \textcolor{matingreen}{فعال} & دوطرفه، با نام فصل و شماره \\
پابرگ & \textcolor{matingreen}{فعال} & با شماره صفحه \\
استایل فصل & \textcolor{matingreen}{فعال} & مقیاس \lr{Boyer/Stewart} \\
فهرست مطالب & \textcolor{matingreen}{فعال} & با نقطه‌چین رنگی \\
پانویس & \textcolor{matingreen}{فعال} & خط \pnum{۳۰٪} عرض \\
فاصله‌ی خطوط & \textcolor{matingreen}{فعال} & \pnum{۱.۱۵} \\
یادداشت حاشیه & \textcolor{matingreen}{فعال} & \lr{\textbackslash mnote} و \lr{\textbackslash marginfig} \\
\end{tblr}
\end{matintable}

\section{نتیجه‌گیری}
\label{sec:summary-conclusion}

\textbf{تمام زیرسیستم‌های صفحه‌آرایی \lr{MatinBook} نسخه‌ی
\lr{\matinbookversion} با موفقیت بارگذاری و آزمایش شدند.
هندسه‌ی صفحه، سربرگ‌ها، پابرگ‌ها، عناوین، فهرست مطالب،
پانویس‌ها، فاصله‌ی خطوط و یادداشت‌های حاشیه‌ای همگی
به‌درستی کار می‌کنند.}

%------------------------------------------------------------------------
% BACK MATTER (symmetric margins)
%------------------------------------------------------------------------
\backmatter
\frontmattergeometry

% Bibliography (empty placeholder)
\printbibliography[title={منابع و مراجع}]

% Index
\printindex

%------------------------------------------------------------------------
% BACK COVER
%------------------------------------------------------------------------
\authorbio{%
    تیم توسعه‌ی \lr{MatinBook}، گروهی از پژوهشگران و برنامه‌نویسان
    فارسی‌زبان است که با هدف ارائه‌ی یک راه‌حل حرفه‌ای برای
    حروف‌چینی کتاب‌های فنی فارسی فعالیت می‌کند. این پروژه حاصل
    سال‌ها تجربه در حوزه‌ی نشر دیجیتال و برنامه‌نویسی است.
}

\makebackcover{%
    این سند، سومین آزمون از مجموعه‌ی آزمون‌های یکپارچگی
    \lr{MatinBook} نسخه‌ی \lr{\matinbookversion} است.

    \vspace{0.5cm}

    \textbf{زیرسیستم‌های آزموده‌شده:}
    \begin{itemize}
        \item هندسه‌ی صفحه
        \item سربرگ و پابرگ دوطرفه
        \item مقیاس عناوین \lr{Boyer/Stewart}
        \item فهرست مطالب رنگی
        \item پانویس و فاصله‌ی خطوط
        \item یادداشت‌های حاشیه‌ای
    \end{itemize}
}

\end{document}